% 原书第六篇6.5节238--248（放射性本节完），印刷259--260，
% 用户原版PDF268--269，书末答案349，PDF358。
\radioactivityproblem{238}{259}
用母原子质量 $M_p$、娘原子质量 $M_d$
和电子质量 $m_e$，
推导 $\beta^-$、$\beta^+$ 与K俘获
放出的反应能 $Q$。
\exsolution
令原子质量包含其相应数目的束缚电子。
$\beta^-$ 衰变原子序数增加1，
核衰变伴随发射一电子，
母/娘中性原子的电子数差恰好抵消发射电子质量，
故
\[
 \boxed{Q_{\beta^-}=(M_p-M_d)c^2}.
\]
$\beta^+$ 衰变原子序数减1，同时发射正电子，
母原子与娘原子电子云计数以及正电子质量
共使静质量差需扣除 $2m_e$：
\[
 \boxed{Q_{\beta^+}=(M_p-M_d-2m_e)c^2}.
\]
K俘获时被俘获的电子已属于母原子，
原子质量的电子数恰能抵消，
\[
 \boxed{Q_{K}=(M_p-M_d)c^2}.
\]
\excheck 各式是中性\emph{原子质量}近似；
如研究俘获后特定壳层的中微子能量，还应扣除
电子空穴的激发能并考虑核反冲。

\radioactivityproblem{239}{259}
用原书原子质量表计算
$\mathrm{^{10}Be}$ 的 $\beta^-$ 衰变最大电子动能，
以及娘核直接处于基态时的反冲动能。
\exsolution
反应为
$\mathrm{^{10}Be}\to\mathrm{^{10}B}+e^-+\bar\nu_e$。
原版质量表中
$\Delta M_{\mathrm{^{10}Be}}\simeq0.01354u$，
$\Delta M_{\mathrm{^{10}B}}\simeq0.01294u$，
故
$Q\simeq(0.00060u)c^2\simeq0.559\,\mathrm{MeV}$。
当中微子能量趋于0时，
电子动能接近最大值，
\[
 \boxed{T_{e,\max}\simeq0.56\,\mathrm{MeV}}.
\]
此时娘核动量与电子相反；
相对论电子满足
$p_e^2c^2=T_e(T_e+2m_ec^2)$，
娘核非相对论反冲能
\[
 \boxed{T_R\simeq
 \frac{T_{e,\max}(T_{e,\max}+2m_ec^2)}
 {2M_dc^2}\simeq47.5\,\mathrm{eV}}.
\]
\excheck 计算反冲必须使用相对论电子动量，
不能简单取 $p_e^2=2m_eT_e$。

\radioactivityproblem{240}{259}
质量 $m=1.0\,\mathrm{mg}$ 的
$\mathrm{^{24}Na}$ 纯 $\beta^-$ 放射性制剂，
半衰期 $T_{1/2}=15$小时。
假设每个 $\beta$ 的平均动能为其端点能的1/3，
求一天内释放的热量。
\exsolution
反应
$\mathrm{^{24}Na}\to\mathrm{^{24}Mg}+e^-+\bar\nu$。
原书质量超额分别为
$-0.00903u$、$-0.01496u$，
所以
$Q=(0.00593u)c^2\simeq5.53\,\mathrm{MeV}$。
按题目假设
$\bar T_\beta=Q/3$。
初核数
$N_0=(10^{-3}\,\mathrm g/24\,\mathrm{g/mol})N_A$，
一天内衰变份额
$f=1-2^{-24/15}$。
故近似放热
\[
 \boxed{W=N_0 f\,\frac Q3
 \simeq5.0\times10^6\,\mathrm J}.
\]
\excheck 题目明确把 $\beta$ 平均动能指定为端点的1/3；
中微子逃逸能不计热量，
娘核反冲及可能随后的其他退激辐射亦忽略。

\radioactivityproblem{241}{259}
用原子质量表计算 $\mathrm{^{11}C}$ 的
$\beta^+$ 衰变中，当娘核\emph{不反冲}时
正电子与中微子的动能。
\exsolution
从原书表得质量超额
$\Delta_{\mathrm{^{11}C}}=0.01143u$、
$\Delta_{\mathrm{^{11}B}}=0.00930u$，
故正电子衰变可用动能
$Q=(0.00213u)c^2-2m_ec^2\simeq0.96\,\mathrm{MeV}$。
娘核零反冲要求正电子与中微子动量相等、方向相反。
令正电子动能 $T$，中微子能量 $E_\nu$，
则
\[
 T+E_\nu=Q,\quad
 E_\nu=pc=\sqrt{T(T+2m_ec^2)}.
\]
两式平方求解
\[
 \boxed{T=\frac{Q^2}{2(Q+m_ec^2)}
 \simeq0.32\,\mathrm{MeV}},
 \quad
 \boxed{E_\nu=Q-T\simeq0.65\,\mathrm{MeV}}.
\]
\excheck 由于正电子有静质量，而中微子近似无质量，
零反冲不意味着两者分得相同能量。

\radioactivityproblem{242}{259}
$\mathrm{^{13}N}$ 的 $\beta^+$ 衰变中，
当正电子动能最大时，
求娘核反冲动能。
\exsolution
衰变为 $\mathrm{^{13}N}\to
\mathrm{^{13}C}+e^++\nu_e$。
原书质量表质量超额
$0.00574u$ 与 $0.00335u$，
故
$Q=(0.00239u)c^2-2m_ec^2
 \simeq1.20\,\mathrm{MeV}$。
最大正电子动能对应中微子能量趋于零，
娘核动量 $p_R\simeq p_e$。
因此
\[
 \boxed{T_R\simeq
 \frac{Q(Q+2m_ec^2)}{2M_Rc^2}
 \simeq0.11\,\mathrm{keV}},
 \quad M_R\simeq13u.
\]
\excheck 这里核反冲能远小于 MeV，
可以在端点近似中令正电子动能 $T\simeq Q$。

\radioactivityproblem{243}{259}
使用原子质量表计算
$\mathrm{^{7}Be}$ 发生 K 电子俘获、
娘核处于基态时的反冲速度。
\exsolution
反应
$\mathrm{^{7}Be}+e^-\to\mathrm{^{7}Li}+\nu_e$。
原书质量超额分别
$0.01693u$ 与 $0.01601u$，
故 $Q\simeq(0.00092u)c^2
 \simeq0.857\,\mathrm{MeV}$。
中微子携带几乎全部释放能量，
动量为 $p_\nu\simeq Q/c$，
娘核反冲速率
\[
 \boxed{v_R\simeq\frac{Q}{M_Rc}
 \simeq3.9\times10^4\,\mathrm{m/s}
 \approx40\,\mathrm{km/s}}.
\]
\excheck K俘获后的壳层空穴可能影响
中微子精细谱线能量；本题要求的数量级可忽略。

\radioactivityproblem{244}{259}
激发 $\mathrm{^{109}Ag}$ 核退激时，
或发射能量 $87\,\mathrm{keV}$ 的 $\gamma$ 光子，
或发射K内转换电子，其结合能为
$26\,\mathrm{keV}$。
求这些电子的速率。
\exsolution
内转换过程将核跃迁能量直接交给壳层电子，
其动能近似
$T_e=87-26=61\,\mathrm{keV}$。
电子总能
$E_e=m_ec^2+T_e=\gamma m_ec^2$，
因此
\[
 \gamma=1+\frac{61}{511},\quad
 \boxed{v_e=c\sqrt{1-\gamma^{-2}}
 \simeq0.45c}.
\]
\excheck 内转换电子并非由 $\gamma$ 光子二次光电吸收产生，
而是核与原子电子直接耦合的退激通道。

\radioactivityproblem{245}{259}
一个初始静止的自由
$\mathrm{^{191}Ir}$ 激发核具有
$E^*=129\,\mathrm{keV}$，
向基态跃迁放出 $\gamma$ 光子。
核反冲引起光子能量的相对改变是多少？
\exsolution
若不计反冲光子能量为 $\varepsilon_0=E^*$。
发射后核动量 $p=\varepsilon/c$，
反冲能
$T_R\simeq\varepsilon^2/(2Mc^2)$。
能量守恒
$E^*=\varepsilon+T_R$。
因 $\varepsilon\ll Mc^2$，
相对光子能量降低为
\[
 \boxed{\frac{\varepsilon_0-\varepsilon}{\varepsilon_0}
 \simeq\frac{E^*}{2Mc^2}
 \simeq3.6\times10^{-7}}.
\]
\excheck 分母中是约 $191u\,c^2$ 的核静能，
不是 $129\,\mathrm{keV}$ 本身。

\radioactivityproblem{246}{259}
自由 $\mathrm{^{191}Ir}$ 核组成的光源与吸收体
发生 $129\,\mathrm{keV}$ 的共振吸收，
为抵消发射、吸收两端的核反冲能差，
两者应有多大相对速率？
\exsolution
自由核发射光子比核能级间隔少
$E_R=\varepsilon^2/(2Mc^2)$；
静止核吸收激发则要求比能级间隔多
$E_R$。两端光子能量失配约
$\Delta E=2E_R=\varepsilon^2/(Mc^2)$。
小速度多普勒能移
$|\Delta\varepsilon|\simeq\varepsilon v/c$，
令其补偿失配，
\[
 \boxed{v=\frac{\varepsilon}{Mc}
 \simeq2.2\times10^2\,\mathrm{m/s}
 =0.22\,\mathrm{km/s}}.
\]
\excheck 这里只求相对速率的模，
必须选择使吸收体观察到光子蓝移的相对运动方向。

\radioactivityproblem{247}{260}
$\gamma$ 射线源位于吸收体上方
$h=20\,\mathrm m$，
为了完全补偿光子向下传播时因地球重力
产生的能量变化，光源应以多大速率向上运动？
\exsolution
弱引力场、低速极限，
光子向下传播的引力能量增量相对为
$\Delta\varepsilon/\varepsilon\simeq gh/c^2$。
光源向上远离吸收体时，
其发射光子相对吸收体发生红移，
一阶相对量约 $v/c$。
令两个效应抵消，
$v/c=gh/c^2$，所以
\[
 \boxed{v=\frac{gh}{c}\simeq
 6.5\times10^{-7}\,\mathrm{m/s}
 =0.65\,\mu\mathrm{m/s}}.
\]
\excheck 问的是光源\emph{向上}运动，
用于抵消沿竖直向下传播光子的引力蓝移。

\radioactivityproblem{248}{260}
$\mathrm{^{67}Zn}$ 激发核发出能量
$\varepsilon=93\,\mathrm{keV}$ 的
Mössbauer 共振 $\gamma$ 光子，
激发态平均寿命 $\tau=14\,\mu\mathrm s$。
源应抬高至少多少，
使在地面观测到的引力谱线位移超过其自然线宽？
\exsolution
弱场引力能量相对变化
$\Delta\varepsilon/\varepsilon\simeq gh/c^2$，
绝对能移
$\Delta\varepsilon\simeq\varepsilon gh/c^2$。
激发态有限寿命导致自然线宽数量级
$\Gamma\simeq\hbar/\tau$。
令 $\Delta\varepsilon\gtrsim\Gamma$，
\[
 \boxed{h_{\min}\simeq
 \frac{\hbar c^2}{g\varepsilon\tau}}.
\]
取 $\varepsilon=93\,\mathrm{keV}$、
$\tau=14\times10^{-6}\,\mathrm s$，
得
$\boxed{h_{\min}\simeq4.6\,\mathrm m}$。
\excheck 使用 $\hbar$ 而非 $h$，
因为自然宽度与角频率寿命关系
$\Gamma\simeq\hbar/\tau$。
